% ED26-REQ-N-c-023-01: Gelfond--Lifschitz 1988, archived author presentation.

\subsection{Gelfond--Lifschitz 1988: stable sets and their autoepistemic translation}
\label{sec:p09-n-c-023-01}
\index{stable model!Gelfond--Lifschitz 1988}
\index{Gelfond--Lifschitz reduct}\index{stable expansion!program translation}

The original construction selects a canonical Herbrand model by testing a
candidate against a program simplified using that candidate.  The paper first
defines \emph{stable sets}, proves that they are minimal models, and then calls
them stable models.  Its phrase \emph{the stable model semantics is defined}
means that there is exactly one such model; the now customary stable-model
relation also covers programs with zero or several stable models.  This
terminological difference does not change the reduct
\cite[\S\S1--2, author-copy pp.~1--5]{GelfondLifschitz1988}.

The source for this treatment is the complete eleven-page archived author
copy, corresponding bibliographically to proceedings pp.~1070--1080, not an
inspected publisher facsimile.  Locators below refer to the author copy's page
sequence and the paper's sections, theorem numbers and program numbers.
Later disjunctive answer sets and complexity classifications are separate
presentations, not results silently attributed to this 1988 paper.

\begin{hbdefinition}{def:p09-gl88-language}{The source's rules, models and grounding}
The source writes a program \(\Pi\) as a set of rules
\[
 A\leftarrow L_1,\ldots,L_m\qquad(m\geq0),
\]
where \(A\) is an atom and each \(L_i\) is an atom or a negated atom.
As an ordinary formula this is
\((L_1\land\cdots\land L_m)\supset A\), universally closed when variables
occur.  In this source-specific notation \(\neg B\) in a rule body denotes
the negative body condition whose candidate-dependent use implements
negation as failure; it is not a strongly negated atom permitted in a head.
When ordinary two-valued modelhood is discussed, that same displayed rule is
read as the indicated classical implication.

The Herbrand universe consists of the language's ground terms, with each term
denoting itself.  A Herbrand interpretation is identified with its set of true
ground atoms.  It is a model when it satisfies every ground rule, and is
\emph{minimal} when no proper subset is a model.  Before the stable-set
construction, replace every rule with variables by \emph{all} its ground
instances.  The resulting program may be infinite; the paper explicitly
allows function symbols and an infinite Herbrand universe.  Write \(At\) for
the set of ground atoms occurring in this grounded \(\Pi\).  Atoms outside
\(At\) cannot be generated by a positive reduct, although a larger Herbrand
base may contain them \cite[\S\S1--2, pp.~1--4]{GelfondLifschitz1988}.
\end{hbdefinition}

For example, source program (2) is
\[
 p(1)\leftarrow,\qquad q(2)\leftarrow,\qquad q(x)\leftarrow p(x).
\]
Over its two constants the Herbrand models are
\(\{p(1),q(1),q(2)\}\) and
\(\{p(1),p(2),q(1),q(2)\}\).  The first is the least model, while the
second contains an unsupported extra fact.  Herbrand selection fixes both
the universe and the interpretations of the constants; arbitrary first-order
models need not do so \cite[\S1, pp.~1--2]{GelfondLifschitz1988}.

\begin{hblemma}{thm:p09-gl88-positive-closure}{proved=annex:p9.n-c-023.1.positive-closure}{Positive programs have a least model and an exact atomic consequence set}
Let \(Q\) be a possibly infinite set of ground rules with finite positive
bodies and one atomic head.  Define
\[
 T_Q(X)=\{A:(A\leftarrow A_1,\ldots,A_k)\in Q,
                    \ \{A_1,\ldots,A_k\}\subseteq X\},
 \qquad X_0=\varnothing,\quad X_{n+1}=T_Q(X_n).
\]
Then \(X_0\subseteq X_1\subseteq\cdots\), and
\(M_Q=\bigcup_{n<\omega}X_n\) is the least Herbrand model of \(Q\).
If \(th(Q)\) means ordinary propositional consequences after treating ground
atoms as propositional letters, then
\[
 th(Q)\cap At=M_Q
\]
for any atom set \(At\) containing the atoms of \(Q\), with absent atoms
false in the least model.  This proves the positive-program dependency used
in the source rather than assuming it; the source invokes that dependency in
\S2 and again in the proof of Theorem~3.
\end{hblemma}

\begin{hbdefinition}{def:p09-gl88-reduct}{The source reduct \(\Pi_M\) and operator \(S_\Pi\)}
For a candidate \(M\subseteq At\), form \(\Pi_M\) by these two operations:
\begin{enumerate}
\item Delete each rule containing a negative body literal \(\neg B\) with
  \(B\in M\).
\item Delete all negative body literals from every remaining rule.
\end{enumerate}
The reduct is positive, so let \(S_\Pi(M)=M_{\Pi_M}\).  The candidate is
a \emph{stable set} precisely when
\[
 S_\Pi(M)=M.
\]
The symbols \(\Pi_M\) and \(S_\Pi\) retain the original presentation;
the modern notation \(P^M\) in the adjacent treatment names the same two
operations, not a different fixed-point procedure.  The source's canonical
model \(CM(\Pi)\) is designated only when this equation has exactly one
solution \cite[\S2, p.~4]{GelfondLifschitz1988}.
\end{hbdefinition}

The candidate makes each body condition \(\neg B\) either unusable
(\(B\in M\)) or dispensable (\(B\notin M\)); positive least-model closure
then checks whether the simplified premises justify exactly that candidate.
In particular, one does not iterate a monotone operator directly on the
original program.  If \(M\subseteq N\), then
\(\Pi_N\subseteq\Pi_M\), hence
\(S_\Pi(N)\subseteq S_\Pi(M)\): the candidate map is antimonotone.
For a positive \(\Pi\), every reduct equals \(\Pi\), so its least model
is its unique stable set \cite[\S\S2--3, pp.~4--5]{GelfondLifschitz1988}.

\begin{hbtheorem}{thm:p09-gl88-theorem1}{proved=annex:p9.n-c-023.1.theorem1}{Original Theorem 1: stable sets are minimal Herbrand models}
Every stable set of \(\Pi\) is a minimal Herbrand model of \(\Pi\).
Consequently stable sets may also be called stable models
\cite[\S2, Theorem~1 and proof, pp.~4--5]{GelfondLifschitz1988}.
\end{hbtheorem}

\begin{hbexample}{def:p09-gl88-one-model}{Source program (5): one nonstratified stable model}
The source's compact program is
\[
 p(1,2)\leftarrow,\qquad q(x)\leftarrow p(x,y),\neg q(y).
\]
Grounding over \(\{1,2\}\) yields the fact and the four rules
\[
\begin{aligned}
 q(1)&\leftarrow p(1,1),\neg q(1),&
 q(1)&\leftarrow p(1,2),\neg q(2),\\
 q(2)&\leftarrow p(2,1),\neg q(1),&
 q(2)&\leftarrow p(2,2),\neg q(2).
\end{aligned}
\]
For \(M=\{q(2)\}\), the reduct contains the fact and
\(q(1)\leftarrow p(1,1)\), \(q(2)\leftarrow p(2,1)\).
Its least model is \(\{p(1,2)\}\), so this candidate is not stable.
For \(M=\{p(1,2),q(1)\}\), the reduct instead contains the fact and
\(q(1)\leftarrow p(1,2)\), \(q(2)\leftarrow p(2,2)\).
It generates exactly \(M\).

Every reduct contains the fact \(p(1,2)\), and none can generate
\(p(1,1),p(2,1),p(2,2)\).  Thus of all \(2^6\) candidates over the
six ground atoms, only four can be fixed points.  Put \(D=\{p(1,2)\}\).
The remaining check is complete:
\[
\begin{array}{c|c}
 M&S_\Pi(M)\\\hline
 D&D\cup\{q(1)\}\\
 D\cup\{q(1)\}&D\cup\{q(1)\}\\
 D\cup\{q(2)\}&D\\
 D\cup\{q(1),q(2)\}&D
\end{array}
\]
The unique stable model is therefore \(D\cup\{q(1)\}\), despite the
program's unstratified negative dependency.  The source reports that its
perfect-model semantics is not applicable and that the well-founded method
selects this same model.  Its separate Prolog observation concerns ground
queries for this particular program, not completeness or termination for
arbitrary Prolog programs \cite[\S3, program~(5), Remarks~1--2,
pp.~5--6]{GelfondLifschitz1988}.
\end{hbexample}

\begin{hbexample}{def:p09-gl88-zero-many}{Source tests with zero or two stable models}
For \(\Pi=\{p\leftarrow\neg p\}\), the only candidates give
\[
 \Pi_{\varnothing}=\{p\leftarrow\},\quad
 S_\Pi(\varnothing)=\{p\};\qquad
 \Pi_{\{p\}}=\varnothing,\quad S_\Pi(\{p\})=\varnothing.
\]
There is no stable model.  Its ordinary classical implication
\(\neg p\supset p\) nevertheless has the minimal model \(\{p\}\),
so minimality alone does not imply stability.

For \(\Pi=\{p\leftarrow\neg q,\ q\leftarrow\neg p\}\), all tests are
\[
\begin{array}{c|c|c}
 M&\Pi_M&S_\Pi(M)\\\hline
 \varnothing&\{p\leftarrow,\ q\leftarrow\}&\{p,q\}\\
 \{p\}&\{p\leftarrow\}&\{p\}\\
 \{q\}&\{q\leftarrow\}&\{q\}\\
 \{p,q\}&\varnothing&\varnothing
\end{array}
\]
Thus exactly \(\{p\}\) and \(\{q\}\) are stable.  Both examples fall
outside the source's unique-canonical-model semantics, though both are
perfectly meaningful tests of the stable-set relation
\cite[\S3, p.~7]{GelfondLifschitz1988}.
\end{hbexample}

\begin{hbcorollary}{thm:p09-gl88-positive-body}{proved=annex:p9.n-c-023.1.positive-body}{The source's empty-model observation}
If every rule of \(\Pi\) has at least one positive body atom, then
\(\varnothing\) is its only stable model.  In particular source program (6),
\[
 p\leftarrow q,\neg r,\qquad
 q\leftarrow r,\neg p,\qquad
 r\leftarrow p,\neg q,
\]
has the unique stable model \(\varnothing\)
\cite[\S3, pp.~6--7]{GelfondLifschitz1988}.
\end{hbcorollary}

Adding the fact \(p(2,1)\) to program (5) changes its behavior.  With
\(D'=\{p(1,2),p(2,1)\}\), the four possible fixed-point candidates give
respectively
\[
\begin{array}{c|c}
 M&S_\Pi(M)\\\hline
 D'&D'\cup\{q(1),q(2)\}\\
 D'\cup\{q(1)\}&D'\cup\{q(1)\}\\
 D'\cup\{q(2)\}&D'\cup\{q(2)\}\\
 D'\cup\{q(1),q(2)\}&D'
\end{array}
\]
so precisely the two singleton choices of \(q\) are stable.  In the source's
game reading, \(p(x,y)\) is a legal move and \(q(x)\) expresses a winning
position: a move to a nonwinning position suffices.  Remarks~3--5 qualify this
reading by the same legal moves for both players and a loop-free game graph.
The added reverse move illustrates the limitation; symmetry alone does not
provide a preferred canonical model.  This is a boundary of that game
interpretation, not a proof that no semantics could ever treat cyclic games
\cite[\S3, pp.~6--7]{GelfondLifschitz1988}.

\paragraph{The original comparison results and their proof boundaries.}
Original Theorem~2 states that if \(\Pi\) has a well-founded model, that
model is its unique stable model.  The paper attributes it to
Van Gelder--Ross--Schlipf [15], Corollary~6.2, rather than supplying a new
proof.  It also attributes strict generality to that source's Examples~6.1--6.2.
In the modern three-valued formulation, the stable-model identification is
the \emph{total} well-founded case, not a claim that every partial
well-founded interpretation is a two-valued stable set; see
Theorem~\ref{thm:p09-well-founded-properties}.
Original Corollary~1 uses [15], Theorem~6.3, to identify the unique stable
model of a locally stratified program with its perfect model.
Original Corollary~2 then identifies the unique stable model of a stratified
program with its iterated-fixed-point model, using the source's cited
perfect-model comparison.  Outside local stratification the source says the
two applicability regions only partly overlap and attributes that observation
to a personal communication.  Its explicit witness in the other direction is
\(p\leftarrow\neg p\), whose sole Herbrand model is perfect in the cited
presentation but is not stable.  These external comparison claims retain
their actual cited standing \cite[\S4, Theorem~2, Corollaries~1--2,
pp.~7--8]{GelfondLifschitz1988}.

\begin{hbdefinition}{def:p09-gl88-expansion}{Source stable expansion and program translation}
The autoepistemic language consists of Boolean formulas and the operator
\(L\); a formula is \emph{objective} if it contains no \(L\).
The source's \(th(X)\) is ordinary propositional consequence, treating each
modal-prefixed formula \(LF\) as a propositional atom in the Boolean
skeleton.  It adds neither modal axioms nor necessitation, as in
Definition~\ref{def:p08-autoepistemic-expansion}.
For an arbitrary set \(E\) of formulas write, as an editorial abbreviation,
\[
 \Delta(E)=\{LF:F\in E\}\cup\{\neg LF:F\notin E\}.
\]
Here \(F\) ranges over \emph{all} formulas, including nested modal formulas,
not merely program atoms.  A stable expansion of \(A\) satisfies
\[
 E=th(A\cup\Delta(E)).
\]
For grounded \(\Pi\), its source translation \(I(\Pi)\) replaces each
negative body literal \(\neg B\) by \(\neg LB\) in the implication
reading of its rule.  Thus
\[
 (A_1\land\cdots\land A_m\land
       \neg B_1\land\cdots\land\neg B_n)\supset H
\]
becomes
\[
 (A_1\land\cdots\land A_m\land
       \neg LB_1\land\cdots\land\neg LB_n)\supset H.
\]
Ground atoms are propositional letters here; \(I\) is a translation into
another logic, not an additional program reduct
\cite[\S4, p.~8]{GelfondLifschitz1988}.
\end{hbdefinition}

\begin{hblemma}{thm:p09-gl88-objective-expansion}{proved=annex:p9.n-c-023.1.objective-expansion}{Objective theories: existence, uniqueness and objective conservativity}
Every objective theory \(A\) has exactly one stable expansion \(E_A\), and
for every objective formula \(F\),
\[
 F\in E_A\quad\Longleftrightarrow\quad A\models F.
\]
The source cites this fact to its references [8] and [6]; the annex supplies
an explicit modal-depth construction and its full proof instead of promoting
the citation to an unproved local premise
\cite[\S4, p.~8]{GelfondLifschitz1988}.
\end{hblemma}

\begin{hblemma}{thm:p09-gl88-expansion-reduct}{proved=annex:p9.n-c-023.1.expansion-reduct}{The source's expansion--reduct lemma}
For every set \(E\) of autoepistemic formulas,
\[
 E\text{ is a stable expansion of }I(\Pi)
 \quad\Longleftrightarrow\quad
 E\text{ is a stable expansion of }\Pi_{E\cap At},
\]
where the reduct on the right is regarded as an objective implication theory
\cite[\S4, lemma and proof, p.~9]{GelfondLifschitz1988}.
\end{hblemma}

\begin{hbtheorem}{thm:p09-gl88-theorem3}{proved=annex:p9.n-c-023.1.theorem3}{Original Theorem 3: unique stable model and unique expansion}
If \(\Pi\) has exactly one stable model \(M\), then \(I(\Pi)\) has
exactly one stable expansion \(E\), and \(M=E\cap At\).
The source expressly credits the lemma-based proof to Halina Przymusinska
\cite[\S4, Theorem~3 and proof, pp.~8--9]{GelfondLifschitz1988}.
\end{hbtheorem}

\begin{hbcorollary}{thm:p09-gl88-correspondence}{proved=annex:p9.n-c-023.1.correspondence}{The same argument yields a correspondence without uniqueness}
The map \(E\mapsto E\cap At\) is a bijection from stable expansions of
\(I(\Pi)\) to stable sets of \(\Pi\).  Its inverse assigns to a stable
set \(M\) the unique expansion of the objective reduct \(\Pi_M\).
This is an explicitly derived consequence of the source's lemma and
objective-theory fact, not the wording of original Theorem~3.
\end{hbcorollary}

An expansion's objective \emph{kernel} is not a complete truth assignment
with every nonderived atom negated.  For example the objective theory
\(A=\{p\}\), in a language also containing \(q\), has an expansion with
\(p,Lp,\neg Lq\), but neither \(q\) nor \(\neg q\).  Thus
\(M=E\cap At\) records derived atoms; it does not identify all objective
formulas in \(E\) with all formulas classically true in the least Herbrand
interpretation.  The objective-expansion proof is needed at the full-formula
level even though the correspondence projects only atoms.

\paragraph{Task and input boundaries of the shared requirement.}
For a finite ground normal input and a fixed candidate, the complete reduct
test deletes the rules, computes positive closure from \(\varnothing\), and
checks equality with the candidate.  At most \(\lvert At\rvert\) strict
stages can occur, so this verification is polynomial in the explicit ground
input size.  Enumerating \(2^{\lvert At\rvert}\) candidates is a finite
decision procedure, not a polynomial enumeration claim.  The original
infinite-ground definition does not inherit that finite bound.
For a fixed function-free rule set with finite input constants, each rule has
a fixed number of variables, so the number of its ground instances is
polynomial in the data size.  When the nonground program itself is input,
that number of variables need not be fixed: a rule with \(k\) independent
variables over \(d\) constants has \(d^k\) instances.  Consequently compact
source size and explicit ground size are distinct measures.

Modern credulous (brave) consequence means truth in \emph{some} stable model,
and skeptical (cautious) consequence means truth in \emph{every} stable model.
In the two-model example, \(p\) is credulous but not skeptical.  In the
source's unique-model case these tests agree.  Under the ordinary universal
convention a no-model program skeptically entails every atom; a nonvacuous
convention would be a different task.  The source's negative body condition
is not the later strong negation of an answer-set literal, and disjunctive
heads change the positive reduct's least-model property to a minimal-model
test.  The shared requirement's later uniform-NP compilation and
normal/disjunctive hardness classifications remain at their separately
attributed treatments in Section~\ref{sec:p09-presentation-n-c-023}; none
is a theorem of the 1988 source.

As an elementary bridge between the modern tasks, adjoining a fresh fact
\(w\leftarrow\) makes every stable model exactly \(M\cup\{w\}\) for an
old stable model \(M\), so existence reduces to credulous \(w\).
Adjoining instead \(w\leftarrow w\) leaves the stable models exactly the
old \(M\), since positive closure never starts that loop, so nonexistence
reduces to skeptical \(w\) under the vacuous convention.  These are explicit
editorial reductions from the reduct definition; they establish no hardness
classification without a separately established hard starting problem.
The same loop is a direct completion boundary: its unique stable model is
\(\varnothing\), whereas the completion \(w\leftrightarrow w\) admits both
classical interpretations.
